% ===================================================================
%  第 8 节 月光与 K3 弦论  原文：TASI10.tex 第 2180--2454 行（PDF p.59--68）
% ===================================================================
\section{月光与 $K3$ 弦论}
\label{chap:moonshineK3}

\begin{keytakeaways}
  \item 核心物理图像：$K3$ 是除四环面外唯一的二维 Calabi--Yau 流形；其 sigma 模型左右两侧的超共形对称性都增强为 ${\cal N}=4$，这正是 Mathieu 月光能在 $K3$ 椭圆亏格中现身的根源。
  \item 数学工具一：把 $K3$ 椭圆亏格按 ${\cal N}=4$ 短/长多重态的特征标展开（式 \eq{massless_characters}--\eq{massive_characters}），有质量多重态的重数 $c_H(n)$ 即 $M_{24}$-模 $K$ 的分次维数。
  \item 数学工具二：改用 ADE 根系标记的奇点椭圆亏格，把同一分解写成 23 个 Niemeier 格对应的等式 \eq{um_splitEG}，从而把 Mathieu 情形推广到全部本影月光。
  \item 重要结论：Witten 指标只收到 $({\rm massless}_L\times{\rm massless}_R)$ 的贡献，故 $K3$ 椭圆亏格由 20 个 ${\cal N}=4$ 短多重态、负系数的短多重态项 $-2\,ch_{1/4,1/2}$ 与有质量项组成；twining 后的 $\phi_g^X$ 只有在 $g$ 保持四平面时才有 $K3$ sigma 模型的几何实现。
  \item 开放问题：只保持两平面的群元素所对应的本影月光函数尚无物理解释；第 \ref{subsec:phys_Other directions} 节列出的非微扰方向（合并对称性、低维紧化、杂化 $K3$、NS5-膜、VOA）目前均无定论。
\end{keytakeaways}

如 \S\ref{sec:Moonshine at weight one-half} 所述，Mathieu 月光最初是在 $K3$ sigma 模型的背景下被注意到的，后来才被认作本影月光 23 个实例之一。下面将会看到，本影月光的全部实例（而不只是 Mathieu 月光）都与同一个量密切相关，即 $K3$ 椭圆亏格。它刻画 $K3$ sigma 模型的 BPS 态，更一般地也刻画含 $K3$ 曲面的超弦紧化的 BPS 态。正是这种与 $K3$ 弦论的可能联系，使 Mathieu 月光与本影月光对许多弦论学家如此有趣。

% [Tutor Note: 此处“$K3$ 椭圆亏格刻画 BPS 态”应理解为：椭圆亏格是保持超对称的指标，只对短多重态（BPS 态）计数，无质量与有质量短多重态的贡献按 ${\cal N}=4$ 特征标分别归类，故它在模空间中不变。]

这是因为 $K3$ 流形有一个特殊性质：它是二维（复维数）唯一的非平凡（即非环面）Calabi--Yau 流形。由于超对称更多，由此得到的弦论紧化比一般 Calabi--Yau 三维流形上的紧化容易处理得多，所含的结构与信息也比环面紧化丰富得多——后者在许多用途上过于简单而无用。因此，$K3$ 紧化在弦论发展史上一直扮演重要角色，是发展革命性新想法的试验场，包括黑洞微观态计数 \cite{StromingerPhys.Lett.B379:99-1041996} 与 AdS/CFT 对应 \cite{MaldacenaAdv.Theor.Math.Phys.2:231-2521998}。此外，它们还出现在联系不同弦论的重要对偶之中。

% [Tutor Note: 黑洞微观态计数指 Strominger--Vafa 1996 年对五维极值黑洞的工作：把 D-膜构型的态计数与 Bekenstein--Hawking 熵比较，是弦论第一次“数出”黑洞熵；$K3\times S^1$ 紧化正是这类构型的基本例子。]

本节将回顾已观测到的月光与 $K3$ sigma 模型之间的关系，并概述为理解这一关系的起源、以及为找到一个能给出人们迫切想要的本影月光模统一构造的物理系统所必须克服的挑战。最后，我们总结针对这些挑战的各种想法及其现状。

\subsection{Mathieu 月光与 $K3$ 椭圆亏格}
\label{sec:M24moonshine}

如 \S\ref{sec:EG} 所述，$K3$ 的 Calabi--Yau 性质连同 Jacobi 形式性质共同决定：它的椭圆亏格由熟悉的弱 Jacobi 形式 \eq{phi01} 给出，即 $2\f_{0,1}$。回忆：除四环面外，$K3$ 是唯一的二维 Calabi--Yau 流形。它有一个特殊性质：不仅 K\"ahler，而且容许超 K\"ahler 结构。因此，第 \ref{sec:EG} 节讨论的 ${\cal N}=2$ 超共形代数的 $U(1)$ 对称性可以扩充为 $SU(2)$，于是 sigma 模型在左、右 mover 两侧都有增强的 ${\cal N}\!=4$ 超共形对称性。这导致 ${\cal N}=4$ 超共形场论的椭圆亏格有一种特定的分解，下面就来解释。

由于底层的共形场论带有 ${\cal N}=4$ 超共形代数的作用，Hilbert 空间分解为该代数（幺正）不可约表示的直和。于是椭圆亏格 (\ref{elliptic_genus}) 可以写成这些表示的特征标之和，每项带某个重数。把 ${\cal N}=2$ 超共形代数的 $U(1)$ 流代数自然地嵌入 ${\cal N}=4$ 超共形代数的 $SU(2)$ 流代数，只须取 $J_0^3 \sim J_0$。因此 ${\cal N}=4$ 的最高权表示仍由两个量子数 $h,\ell$ 标记，分别对应算子 $L_0$、$J_0^3$；这样的不可约表示 $V_{h,\ell}$ 的特征标定义为
$$
{\it ch}_{h,\ell}(\t,z) := \tr_{V_{h,\ell}} \left( (-1)^{J_0}y^{J_0} q^{L_0-c/24}\right) .
$$
% [Tutor Note: 这里 $J_0^3\sim J_0$ 只表示二者成比例，比例常数取决于 $U(1)$ 流在 $SU(2)$ 中的嵌入约定（常见约定是 $J^{{\cal N}=2}=2J_0^3$）。记号 $\sim$ 正是为了回避这个常数：一旦取定归一化，特征标里 $y^{J_0}$ 的荷与 $\ell$ 的对应才确定，草率地把 $J_0^3$ 当成 $J_0$ 会让 $y$ 的幂次差一个因子。]
对中心荷 $c=6$，Ramond 扇区有两个超对称表示（也称“BPS”表示或“无质量”表示），其量子数为
$$
h=\tfrac{1}{4}\;,\; \ell = 0 ,\tfrac{1}{2} \;.
$$
它们的特征标已算出 \cite{Eguchi1987,Eguchi1988,Eguchi1988a}，为
\bea\label{massless_characters}
{\it ch}_{\frac{1}{4},0}(\t,z) &=& \frac{\th_1^2(\t,z)}{\eta^3(\t)} \m(\t,z)\, ,\\\notag
{\it ch}_{\frac{1}{4},\frac{1}{2}}(\t,z) &=& q^{-\frac{1}{8}} \frac{\th_1(\t,z)}{\eta^3(\t)} - 2 \frac{\th_1(\t,z)^2}{\eta^3(\t)} \m(\t,z)\, ,
\eea 
其中 $\m(\t,z)$ 是所谓的 Appell--Lerch 和，定义见 \eq{def:mu}。

% [Tutor Note: $\m$ 是 Appell--Lerch 和，属于模拟 Jacobi 形式：它本身不满足 Jacobi 形式的模变换，而带一个由影子决定的非全纯修正。正因如此，(8.1) 的两个特征标单独看都不是弱 Jacobi 形式；只有按 (8.3) 组合、并让 $\m$ 与 $\theta_1^2/\eta^3$ 的因子相互补齐之后，$K3$ 椭圆亏格才是模不变的。忽略这一点会把“模拟”误当成“全纯”。]

注意，超对称表示的 Witten 指标不为零：
% [Tutor Note: 原文 p.61 此处第二个式子仍写作 $ch_{1/4,0}$，按上下文（$h=1/4$ 的短多重态有 $\ell=0$ 与 $\ell=1/2$ 两个）应为 $ch_{1/4,1/2}(z=0)=-2$。译文照录原文，读者勿从字面读作两个相同之量。]
$$ch_{\frac{1}{4},0}(\t,z=0) =1\,, \quad ch_{\frac{1}{4},0}(\t,z=0) =-2\,.$$ 
非超对称表示（也称“非 BPS”表示或“”）具有
% [Tutor Note: 原文 p.61 此句第三个引号内空缺，按术语表与上下文应为 massive（“有质量”）。此处照录空缺，语义即“非超对称（或称非 BPS、有质量）表示”。]
$$
h =\tfrac{1}{4} + n\, ,\quad \ell =\tfrac{1}{2}\, ,\quad\text{} \quad n = 1,2,\dots,
$$
其特征标为
\be\label{massive_characters}
ch_{\frac{1}{4}+n,\frac{1}{2}}(\t,z) = q^{-\frac{1}{8}+n} \frac{\th_1^2(\t,z)}{\eta^3(\t)} \, .
\ee
% [Tutor Note: 幂次 $q^{-1/8+n}$ 极易抄错：$c=6$ 时 $-c/24=-1/4$，但 (8.1) 中的指数是 $-1/8$（$c/24$ 的 $1/4$ 已由 $\theta_1,\eta$ 的 $q^{1/8}$ 因子吸收），不要直接写成 $-1/4$。每升高一档有质量能级就多一个 $q$，故一般项为 $q^{-1/8+n}$；$n$ 从 $1$ 起，$n=0$ 的边界即短多重态 $h=1/4$，不属有质量谱。]
按定义，它们对 Witten 指标的贡献为零，即 $ch_{\frac{1}{4}+n,\frac{1}{2}}(\t,z=0) = 0$。这就是为什么 Witten 指标只收到 ${\cal N}=4$ 表示中 (massless$_L\times{\rm massless}_R$) 型量子态的贡献，而椭圆亏格既收到 (massless$_L\times{\rm massless}_R$) 的贡献，也收到 (massive$_L\times{\rm massless}_R$) 的贡献。

现在可以给出 $K3$ 椭圆亏格按 (\ref{splitEGM24}) 分解的第一种物理解释，这一解释使与 $M_{24}$ 的联系变得明显。由 (\ref{splitEGM24}) 以及 \eq{massless_characters}--\eq{massive_characters} 可见，${\bf EG}(K3)$ 可表示为
\be
{\bf EG}(K3) =  20\, ch_{\frac{1}{4},0} -2 \, ch_{\frac{1}{4},{1\over 2}}  + \sum_{n=1}^\infty c_H(n)\,ch_{\frac{1}{4}+n,\frac{1}{2}}
\ee
其中
\be 
H(\t) =q^{-1/8}\left(-2 +\sum_{n=1}^\infty c_H(n)q^n\right) =2 q^{-\frac{1}{8}} \left( -1 + 45\, q+ 231\, q^2 + 770\,q^3\dots \right)
\ee
是 \eq{second_def_H} 给出的模拟模形式。特别地，按 Mathieu 月光猜想 \ref{conj_mock}，有质量多重态的重数 $c_H(n)$ 等于 $M_{24}$-模 $K$ 的齐次分量的维数。

% [Tutor Note: 由 \eq{massive_characters} 后的展开可见 $c_H(1)=90$、$c_H(2)=462$、$c_H(3)=1540$，正是 $M_{24}$ 不可约表示的维数组合：$90=45+45$、$462=231+231$、$1540=770+770$。式中的 $45,231,770$ 是 $H$ 括号内的系数，因整体乘以 $2$ 才得到 $c_H(n)$；漏掉这个因子是常见错误。]

如 \S\ref{subsec:Mathieu} 已提到的，于是人们很自然地想把 $\f_g$ \eq{K3_EG_Fou}（它编码 Mathieu 月光的 McKay--Thompson 级数 $H_g$）认作 $K3$ 非线性 sigma 模型在某个保持 ${\cal N}=4$ 超共形代数的对称性下的 twined 椭圆亏格。对某些 $g\in M_{24}$ 这确实可行，但对另一些群元素，这一解释已被证明不成立。我们将在 \S \ref{sec:symK3} 讨论这一点。

\subsection{本影月光与 $K3$ 椭圆亏格}
\label{sec:umK3}

如前所暗示，对椭圆亏格的分解 (\ref{splitEGM24})——一部分直接来自月光的模拟模形式，另一部分来自 Appell--Lerch 和 $\mu$——${\cal N}=4$ 分解并不是唯一可能的解释。第二种解释以 ADE 根系标记的奇点空间的椭圆亏格为语言，在本例中给出相同结果，但关键差别在于：它可以立刻推广到本影月光其余 22 个实例。

% [Tutor Note: “其余 22 个”来自 Niemeier 格的分类：秩 24 的偶自对偶格共 24 个，其中 Leech 格无根，剩下 23 个带根系的 Niemeier 格 $N^X$ 对应 23 个本影实例；$X=24A_1$ 即 Mathieu 情形，故“其余 22 个”是分类的直接结果，不是巧合数字。]

为说明这一点，我们先给 23 个根系 $X$ 各配一个函数 $\eg(X;\tau,z)$。注意 $K3$ 曲面的一种常见定义是：它至多有 du Val 型曲面奇点，即复平面上形如 $\mathbb C^2/G$ 的奇点，其中 $G$ 是 $SU(2)_{\mathbb C}$ 的有限子群。按著名的 McKay 对应，这些奇点有 ADE 分类。以这些 ADE 奇点为靶空间时，弦论的共形场论描述由 \cite{Ooguri:1995wj} 给出。其椭圆亏格的形式已在多篇文章中研究过，包括 \cite{Troost:2010ud,Eguchi:2010cb,Ashok:2011cy,Ashok:2013pya,Murthy:2013mya,Cheng:2014zpa,Harvey:2014nha}。记 $\eg(Y;\tau,z)$ 为对应于根系 $Y$ 的奇点的椭圆亏格的全纯部分（另见 \S\ref{sec:EG} 末的注记）。把定义扩充到根系之并，$\eg(\cup_i Y_i;\tau,z)= \sum_i \eg( Y_i;\tau,z)$，就得到 23 个 Niemeier 根系 $X$ 各自的 $\eg(X;\tau,z)$ 的定义。例如
\be
\eg(A_1;\tau,z) = ch_{\frac{1}{4},0}(\t,z) = \frac{\th_1^2(\t,z)}{\eta^3(\t)} \m(\t,z), 
\ee
而 $\eg(24A_1;\tau,z) = 24\frac{\th_1^2(\t,z)}{\eta^3(\t)} \m(\t,z)$ 正是 \eq{splitEGM24} 中 $\eg(K3)$ 分解的第一项。

% [Tutor Note: 单个 $A_1$ 奇点给出 $ch_{1/4,0}$ 而非 $ch_{1/4,1/2}$：$A_1$ 是 $\CC^2/\ZZ_2$ 型 du Val 奇点，其椭圆亏格只含无质量表示中最简单的一支。系数 $24$ 来自 $24A_1$ 中 24 个分量相加，故与 (8.3) 中 $20\,ch_{1/4,0}-2\,ch_{1/4,1/2}$ 的 $20$ 不是同一类计数，勿混。]

回忆：对 23 个 Niemeier 格 $N^X$（由根系 $X$ 唯一确定），本影月光对每个 $N^X$ 各配一个有限群 $G^X$，并对每个 $g\in G^X$ 配一个权 1 的模拟 Jacobi 形式 $\psi^X_g = \sum_r H^X_{g,r} \tilde\theta_{m,r}$。\cite{Cheng:2014zpa} 证明 $\eg(K3)$ 有 23 种分成两部分的方式：
\be\label{um_splitEG}
{\bf EG}(K3;\tau,z) = \eg(X;\tau,z) + {\theta_1^2(\tau,z) \over 2\eta^6(\tau)} \left({1\over 2\pi i}{\partial\over \partial w} \psi^X(\tau,w)\right) \Big\lvert_{w=0}
\ee 
其中 $\psi^X =\psi^X_e$ 是 \S\ref{sec:umbral} 给出的最优模拟 Jacobi 形式，编码本影月光模的分次维数；第一项就是上面讨论的奇点椭圆亏格。
% [Tutor Note: 第二项先对 $w$ 求导再取 $w=0$，且 $\theta_1^2$ 前的系数是 $1/(2\eta^6)$。抄写时最易漏两点：分母上的 $2$ 与导数前的 $1/(2\pi i)$。漏掉前者会与 (8.7) 的 $1/\eta^6$ 混淆（见 (8.7) 前批注），漏掉后者会使整项差一个 $2\pi i$ 因子，破坏与 (8.3) 的逐项对照。]
因此，椭圆亏格的分解 \eq{splitEGM24} 也可解释为两项贡献之和：一项来自根系 $X=24A_1$ 对应的奇点椭圆亏格，另一项来自本影月光函数 $\psi^{24A_1}$。与 ${\cal N}=4$ 特征标解释不同，这种对分解的解释适用于本影月光全部 23 个实例。换言之，\eq{um_splitEG} 实际上是 23 个等式，对应全部 23 个 Niemeier 格 $N^X$。

这一解释提示函数 ${\bf EG}(K3;\tau,z)$ 有一个 twined 版本：把等式 \eq{um_splitEG} 的右端 twining 即可：
% [Tutor Note: 注意 (8.6) 的 $\theta_1^2$ 项系数是 $1/(2\eta^6)$，而 (8.7) 写成 $1/\eta^6$；取 $g=e$ 并对 (8.6) 右端逐项 twining，两式理应给出同一个 $\phi_e^X={\bf EG}(K3)$，故这里相差因子 2，疑为原文笔误。译文照录，定稿前宜按 CDH 本影月光原文核对。]
\be\label{um_twining}
\phi_g^{X}(\tau,z) := \eg_g(X;\tau,z) + {\theta_1^2(\tau,z) \over \eta^6(\tau)} \left({1\over 2\pi i}{\partial\over \partial w} \psi^X_g(\tau,w)\right) \Big\lvert_{w=0}~,
\ee 
其中 $g$ 可取 $G^X$ 中任意元素。右端第一项是对应奇点 sigma 模型的 twined 椭圆亏格（的全纯部分）。它可以显式计算：我们已经知道 $G^X$ 在 Niemeier 根系 $X$ 上的作用，而后者翻译为奇点共形场论上保持超共形结构的 $G^X$ 作用，由此可以定义其 twined 椭圆亏格。第二项则由本影月光的 McKay--Thompson 级数 $\psi^X_g$ 直接给出。

% [Tutor Note: 这里的 twining 是逐项进行的：第一项 twining 奇点共形场论，第二项 twining 本影 McKay--Thompson 级数，并没有在 $K3$ 的 Hilbert 空间上取迹。因此 $\phi_g^X$ 只是由 (8.7) 形式定义出的函数；它是否真能写成某个 $K3$ sigma 模型对称性的 twined 椭圆亏格，正是 8.4 节要检验的问题。]

容易验证：取 $N=N^{24A_1}$ 时，$\phi_g^X$ 的这一定义给出与 \eq{K3_EG_Fou} 中 Mathieu 月光情形相同的弱 Jacobi 形式。

% [Tutor Note: 作者略去了验证细节：$N=24A_1$ 时本影模拟 Jacobi 形式 $\psi^{24A_1}_g$ 按 $\tilde\theta_{1,r}$ 分解后的系数，恰好复现 Mathieu 月光的 $H_g$；这一步依赖 24A_1 的 $G^X\cong M_{24}$ 以及指标 $m=1$ 的 theta 分解，参见 Cheng--Duncan--Harvey 的本影月光构造。]

与 Mathieu 月光的情形一样，人们很自然地想把由 \eq{um_twining} 得到的全部本影月光函数 $\phi_g^{X}$ 都解释为 $K3$ 非线性 sigma 模型在某个保持 ${\cal N}=4$ 超共形代数的对称性下的 twined 椭圆亏格。对某些 $g\in G^X$，这样的解释并不存在，详见 \S\ref{sec:symK3}。具体地说，只有 $G^X$ 中保持一个四平面（即 $N^X\otimes_{\ZZ} \RR$ 内的四维定向子空间）的子群，才能充当 $K3$ sigma 模型的对称群。为后文方便，对给定带非平凡根系的 Niemeier 格 $N=N^X$，把 \eq{um_twining} 给出的 Jacobi 形式之集记作
\be\label{def:umset}
\Phi(N^X):= \{\phi_g^X~\lvert~g~{\text{为 $G^X$ 的保持四平面的元素}}\}~.\ee

\subsection{Conway 月光与 $K3$ 椭圆亏格}
\label{subsec:conwayK3}

回忆 \S \ref{sec:conway_moonshine} 描述的 Conway 扭模（twisted Conway module）$V_{\text{tw}}^{s\natural}$ 的构造。给定 $\Lambda \otimes_{\ZZ} \CC$ 中一个固定的 $n$ 维子空间，把 $24$ 个费米子作双线性组合，可以造出不同的 $U(1)$ 流 \cite{Cheng:2014owa,MR3465528}。在这个理论中，与一个四平面（即实维数为四的子空间）相关联的费米子可以造出水平为 2 的 $U(1)$ 流 $J$。给定这样的四平面 $\Pi$，固定 $U(1)$ 流与 SVOA 结构，会把 $V^{s\natural}_{\text{tw}}$ 的对称性从 $\Co_0$ 破缺到其保持 $\Pi$ 的子群 $G_\Pi$。反之，给定保持四平面的 $G_\Pi< \Co_0$，可以造出 $U(1)$ 流 $J$，使得 $V^{s\natural}_{\text{tw}}$ 在赋予 $J$ 与 ${\cal N}=1$ 超共形代数的模结构之后，对称群恰为 $G_\Pi$。

有趣的是，$V^{s\natural}_{\text{tw}}$ 的带 $U(1)$ 荷的分次特征标与 ${\bf EG}(K3)$ 一致（相差一个符号）\cite{MR3465528}。更一般地，可以考虑 Conway 扭模在保持四平面的 $\Co_0$ 元素下 twining 所给出的 $U(1)$-分次特征标（记号见 \S \ref{sec:conway_moonshine}）：
\begin{equation}
\f_g^\Lambda:=-\tr_{V^{s\natural}_{\text{tw}}}\left[\mathfrak{z}~\hat{g}~y^{J_0}q^{L_0-\frac{c}{24}}\right],
\end{equation}
其中 $J_0$ 是 $U(1)$ 流的零模。显式地，它由下式给出：
\begin{equation}
\label{eq:phi_g}
\begin{split}
\f_g^\Lambda(\tau,\zeta)&=\frac{1}{2}\left[\frac{\theta_3(\tau,\zeta)^2}{\theta_3(\tau,0)^2}\frac{\eta_{-g}(\tau/2)}{\eta_{-g}(\tau)}-\frac{\theta_4(\tau,\zeta)^2}{\theta_4(\tau,0)^2}\frac{\eta_{g}(\tau/2)}{\eta_{g}(\tau)}\right. \\
&~~~~~~~~~~\left.-\frac{\theta_2(\tau,\zeta)^2}{\theta_2(\tau,0)^2}C_{-g}\eta_{-g}(\tau)-\frac{\theta_1(\tau,\zeta)^2}{\eta(\tau)^6}D_g\eta_g(\tau)\right]~,
\end{split}
\end{equation}
其中大部分特殊函数已在 \eqref{eta_g_and_C_g} 中定义。函数 $D_g$ 定义为
\begin{equation}
D_g:=\n\sideset{}{'}\prod_{i=1}^{12}\left(1-\l_i^{-1}\right)=\sideset{}{'}\prod_{i=1}^{12}\left(\n_i-\n_i^{-1}\right).
\end{equation}
其中在 $\sideset{}{'}\prod$ 里跳过与固定四平面对应的两对特征值，对其有 $\l_i^{\pm1}=1$。注意 $D_g$ 非零，当且仅当 $g$ 恰好固定一个四平面而不固定更多。在后一种情形（恰好固定一个四平面），$g$ 的特征值把 $D_g$ 确定到相差一个符号，因为我们可以自由交换 $\l_i$ 与 $\l_i^{-1}$ 的称呼。因此，对恰好保持四平面的元素，$\f_g^\Lambda$ 其实有两种取法，取决于 $D_{g}$ 的符号；于是定义
\begin{equation}
\begin{split}\label{eq:phi_epsilon}
\phi^\Lambda_{g,\e}(\tau,\zeta)&:=\frac{1}{2}\left[\frac{\theta_3(\tau,\zeta)^2}{\theta_3(\tau,0)^2}\frac{\eta_{-g}(\tau/2)}{\eta_{-g}(\tau)}-\frac{\theta_4(\tau,\zeta)^2}{\theta_4(\tau,0)^2}\frac{\eta_{g}(\tau/2)}{\eta_{g}(\tau)}\right. \\
&~~~~~~~~~~\left.-\frac{\theta_2(\tau,\zeta)^2}{\theta_2(\tau,0)^2}C_{-g}\eta_{-g}(\tau)-\e~\frac{\theta_1(\tau,\zeta)^2}{\eta(\tau)^6}\left|D_g\right|\eta_g(\tau)\right]~,
\end{split}
\end{equation}
其中 $\e=\pm1$ 编码 $D_g$ 的符号歧义。

% [Tutor Note: (8.10) 与 (8.12) 中 $\theta_3$ 项配 $\eta_{-g}$、$\theta_4$ 项配 $\eta_{g}$，下标的正负号容易抄反；一旦把 $g$ 换成 $-g$，$\theta_3$ 与 $\theta_4$ 两项互换，而 $\theta_2$ 项不变。另注意 $D_g$ 只出现在 $\theta_1^2/\eta^6$ 项，且 (8.12) 用 $|D_g|$ 代替 $D_g$。]

在所有这些情形，\cite{MR3465528} 证明：对每个至少固定一个四平面的 $g\in\Co_0$，$\f_{g,\e}^\Lambda$ 都是权 $0$、指标 $1$ 的 Jacobi 形式，只是模群要取某个 $\Gamma_g\subseteq \SL(2,\ZZ)$。为后文方便，把所有这些函数之集记作
\begin{equation}\label{def:conwayEG}
\Phi(\Lambda):=\left\lbrace \phi_{g,+}^\Lambda, \phi_{g,-}^\Lambda~|~g~\text{为 $\Co_0$ 的保持四平面的元素}\right\rbrace~.
\end{equation}

% [Tutor Note: 这里的结论不是对一般 $\Co_0$ 元素成立，而是逐个共轭类验算的结果（$g$ 与 $-g$ 的 $\eta$ 商配对给出乘子系统）。验证要点是 $\Gamma_g$ 取为稳定相应四平面与特征值的同余子群，模权为 0、指标为 1；$\Gamma_g$ 一般不是 $\SL(2,\ZZ)$ 本身，这一点在读“Jacobi 形式”时勿默认全模群。]

关于这组函数与本影月光给出的那组函数（即 \eq{def:umset} 定义的 $\Phi(N^X)$），我们做几点观察。其一，这两组函数明显有大面积重叠：Conway 月光与本影月光常常给出同一个弱 Jacobi 形式，尽管途径大不相同。一个重要差别是：通过本影月光 \eq{um_twining}，弱 Jacobi 形式 $\phi^X_g$ 对\underline{所有} $g\in G^X$ 都有定义；而在 Conway 情形，只能限于保持四平面的元素（这样才能得到指标为 1 的 Jacobi 形式）。其二，只出现在 $\Phi(\Lambda)$ 中的函数，往往实现为 $K3$ sigma 模型的 twined 椭圆亏格，位置在模空间里那些模型可描述为 $T^4$-轨道流形的子轨迹上。也请注意，$T^4$-轨道流形的几何与 Conway 模的自由场轨道流形构造衔接得很好。事实上，当人们（粗略地说）“忘记”左右 mover 的区分时，$V_{\text{tw}}^{s\natural}$ 同构于某个 $T^4$ 模型的 $\Z_2$ 轨道流形 \cite{MR3465528,Taormina:2017zlm,Creutzig:2017fuk}。其三，只来自本影月光、不出现在 Conway 月光中的函数，必定对应一种在 $K3$ sigma 模型左右 mover 上作用不同的对称性 \cite{ChengMirandaC.N.2016KSTL}。下一小节将讨论月光对称群与弦论 $K3$ 对称群之间的联系。

% [Tutor Note: 第二点中“忘记左右 mover 的区别”是技术性说法：Conway 模是手征的（只含一种手征性），而 $T^4$ 模型有左、右两支；做 $\Z_2$ 轨道流形并取手征一半，才与 $V^{s\natural}_{\text{tw}}$ 同构。因此这条对应只适用于 $\Phi(\Lambda)$ 中由 $T^4$-轨道流形子轨迹实现的那部分函数。]

\subsection{$K3$ 弦论的对称性与月光}
\label{sec:symK3}

前几个小节展示了 $K3$ 椭圆亏格与本影、Conway 月光函数之间密切的数学关系。因此，比较月光群与 $K3$ sigma 模型的对称性是很自然的。事实上，人们早已知道 $K3$ 曲面的对称性与散在群有密切关系。特别地，Mukai 的一个著名定理 \cite{Mukai} 确立了 $M_{23}$ 的角色：这个 26 个散在群之一、同时是 $M_{24}$ 子群的群，正是“组织”$K3$ 曲面对称性的有限群。更精确地说，\cite{Mukai} 证明：至少带五个轨道的 $M_{23}$ 子群（的同构类）与 $K3$ 辛同构的有限群（的同构类）之间存在一一对应。

% [Tutor Note: “至少五个轨道”是定理的精确条件（$M_{23}$ 作用在 24 个点上）；文献中也有写作“五个轨道”的表述，原文取“至少”，译文照录。该条件不可省去，否则对应关系不成立；证明见 Mukai 原文与 Kondo 的简化证明。]

Gaberdiel、Hohenegger 与 Volpato 启动了这一经典结果向“弦论 $K3$ 几何”的推广：他们在 \cite{Gaberdiel2011} 中用格论技巧分类了所有（非奇异）$K3$ sigma 模型的对称群，方法紧接 Kondo \cite{Kondo} 对 Mukai 定理的简化证明。从时空（D-膜）视角看，\cite{Gaberdiel2011} 的结果可以看作分类 $K3$ 曲面上导出范畴的辛自等价（对称性）\cite{MR3586514}。关于与 $K3$ 上曲线计数相关的适当模空间的对称性，相关讨论亦可参见 \cite{Cheng:2015kha}。

\cite{ChengMirandaC.N.2016KSTL} 把分类推广到对应于 $K3$ sigma 模型模空间中奇点轨迹的理论。如果关心的是完整弦论而非仅有 sigma 模型，这一步是必需的。尽管 IIA 型世界面理论在这些轨迹上行为很差，完整的 IIA 型弦论不仅完全定义良好，而且在非阿贝尔规范对称性方面具有特殊的物理意义 \cite{Aspinwall:1995zi}。在这些轨迹上，时空规范群从 $U(1)^{24}$ 增强为某个非阿贝尔群，其 ADE 型规范群由 $K3$ 曲面的 ADE 型奇点给出 \cite{Witten:1995ex,Aspinwall:1995zi}。模空间中存在这类规范对称性增强的轨迹，虽然从 IIA 型世界面分析并不直接显现，但从杂化 $T^4$ 紧化的对偶描述却很清楚。

% [Tutor Note: 这段的物理图像是：在小体积极限下世界面 sigma 模型退化（世界面理论“行为很差”），但完整弦论仍良好定义；规范群增强由 $K3$ 上收缩的二维球（ADE 奇点的例外曲线）给出。原文强调“对偶描述很清楚”，指的是 IIA 在 $K3$ 上的紧化与杂化弦在 $T^4$ 上的紧化互为对偶，后者在增强点上没有世界面奇异性。]

为陈述分类结果，先回顾 K3 上 sigma 模型的一些一般性质（见 \cite{Aspinwall:1996mn,Nahm:1999ps}）。具有 $\mathcal{N}=(4,4)$ 超对称的非奇异非线性 sigma 模型在 $K3$ 上的模空间，是下面空间中的一个开子集：
\be \label{moduli_space}{\cal M}=(SO(4)\times O(20))\backslash O^+(4,20)/O^+(\Gamma^{4,20})\ ,
\ee
其中 $(SO(4)\times O(20))\backslash O^+(4,20)$ 是正四平面 $\Pi$ 的 Grassmannian，即 $\RR^{4,20}\cong \Gamma^{4,20}\otimes_\ZZ \RR$ 内的四维定向正定子空间；$\Gamma^{4,20}$ 是符号差为 $(4,20)$ 的偶幺模格。换言之，在模空间 ${\cal M}$ 中取一点，等价于在 $\Gamma^{4,20}\otimes_\ZZ \RR$ 中取一个四平面。这一开子集在 $\M$ 中的补集由奇异四平面 $\Pi$ 组成（即与某个根 $v\in \Gamma^{4,20}$，$v^2=-2$ 正交），它们对应 K3 上 sigma 模型的某些奇异极限。整个空间 ${\cal M}$ 也是固定有限 $g_s$ 值下 IIA 型弦论的模空间。

在奇异情形，相应 IIA 型弦论的全部保超对称对称性含有一个连续非阿贝尔规范对称性；把它模掉，就得到离散对称群 $G_{\rm IIA}$。于是
\be
G_{\rm IIA}(\Pi) = {\rm Stab}(\Pi)/{\rm W} 
\ee
其中 ${\rm Stab}(\Pi)$ 是 $O^+(\Gamma^{4,20})$ 中逐点稳定四平面 $\Pi$ 的子群，${\rm W}$ 是由关于所有与 $\Pi$ 正交的根向量（若存在）的反射生成的 Weyl 群。

% [Tutor Note: (8.14) 是双陪集，不是直积：$O^+(4,20)$ 中左右分别模去 $SO(4)\times O(20)$ 与 $\Gamma^{4,20}$ 的算术子群。符号差写作 $(4,20)$ 而非 $(20,4)$，正定四平面 $\Pi$ 是“正”的那一半；若把约定反过来，根 $v^2=-2$ 会变成 $+2$，8.5 节讲奇异轨迹时就会出错。]

随后 \cite{ChengMirandaC.N.2016KSTL}（推论 4）证明：在 ${\cal M}$ 中某处实现的群 $G_{\rm IIA}$，与 24 个秩 24 偶自对偶格所对应的 24 个有限群 $G_N$（参看 \eq{def:niemeier_grp}）的保持四平面的子群一一对应；这 24 个有限群包括 23 个本影群 $G^X$ 与 Conway 群 $\Co_0$。特别地，这意味着一般而言本影群 $G^X$——尤其是 $G^{24A_1}\cong M_{24}$——{\underline{不可能}}是任何单个 $K3$ sigma 模型的对称群，因为它们不是 $\Co_0$ 的保持四平面的子群。

% [Tutor Note: 这一“不可能”是本节的转折点：$M_{24}$ 含有大量只保持两平面的元素，故它无法作为单个 $K3$ sigma 模型的对称群出现。这正是 8.5 节必须离开微扰 sigma 模型框架的原因，也是本影月光“物理起源”问题的核心难点。]

把对称群作为抽象群分类之后，我们还想知道它们如何作用在超对称态的谱上。这一信息由这些对称性的 twined 椭圆亏格承载，其定义见 \eq{def:twined_elliptic_genus}。\cite{ChengMirandaC.N.2016KSTL} 猜想（连同其他更精细的条件）：在模空间 ${\M}$ 中任何位置实现的 twined $K3$ 椭圆亏格，都至少与 \eq{def:umset} 和 \eq{def:conwayEG} 定义的 24 组函数 $\Phi(N)$ 中某一组里的某个函数一致；这里 $N$ 或是 23 个 Niemeier 格 $N^X$ 之一，或是 Leech 格 $\Lambda$。反之亦猜想：$\Phi(N)$ 的{\it 所有}元素都在刻画 $K3$ 的 BPS 态对称性中起作用。

% [Tutor Note: 这条猜想有两半：前半说“所有 twined $K3$ 椭圆亏格都落在 $\Phi(N)$ 里”（满射方向），后半说“$\Phi(N)$ 的每个元素都被某个物理对称性实现”（单射方向）。8.5 节开头说的“不足以解释本影月光”，针对的是后半：保持两平面的元素对应的 $\phi^X_g$ 不在 $\Phi(N)$ 的物理可实现部分中。]

最近，\cite{Paquette:2017gmb} 考察了与紧化在 $K3 \times T^2$ 上的 IIA 型弦论密切相关的弦论（通过取轨道流形），并要求 BPS 计数函数的解析结构与壁穿越物理相容，从而为上述猜想找到了有力证据。然而，人们还不理解为什么本影月光函数会与 $K3$ sigma 模型的 twined 椭圆亏格有如此密切的关系。更多讨论见 \S\ref{subsec:phys_Other directions}。

% [Tutor Note: 此处“wall-crossing（壁穿越）”指 BPS 态计数在模空间中的跳变：跨过边际稳定壁时 BPS 谱改变，但指标（含 twining 的指标）在适当修正下不变。Paquette 的工作正是用这一约束把可实现函数集合压到 $\Phi(N)$ 上。]

\subsection{超出微扰弦论}
\label{subsec:phys_Other directions}

上一小节最重要的启示是：$K3$ 共形场论的保超对称对称性与本影月光的对称性密切相关，但不足以解释本影月光——许多本影群根本无法在模空间的任何位置实现为 $K3$ sigma 模型的对称性。

具体地说，\S\ref{sec:symK3} 讨论的猜想断言：与保持四平面的群元素相关联的全部本影、Conway 月光函数，都作为 $K3$ sigma 模型的 twining 亏格起作用，并且事实上完全刻画了后者；然而，只保持两平面的群元素所对应的本影（含 Mathieu）月光函数有什么物理意义，仍不清楚。

% [Tutor Note: 本段是全文的“承认未知”处：作者明确说缺乏共识。注意区分两个层次——(1) 保持四平面的子群有 $K3$ sigma 模型实现；(2) 只保持两平面的元素（例如 $M_{24}$ 中大量元素）没有已知的几何实现。本影月光的 23 个实例中，正是这些“两平面”元素使问题超出微扰弦论。]

关于什么是与本影月光相关的确切物理设定，目前尚无共识。与此同时，恰恰是这种超出微扰弦论的必要性，以及让月光进入弦论黑洞等非微扰设定、或进入某种新型超对称结构的可能性，使本影月光与弦论之间可能的联系如此令人兴奋。下面简要讨论文献中已有不同程度探讨的几种可能。

\myparagraph{合并对称性。} 其想法是找到一种办法，把在模空间不同点实现的对称性合并起来，从而生成一个更大的群，一个不受限于模空间某个特定点的“meta”群；它还要包含保持两平面的元素。

这一可能性最早以问题的形式出现在原始文章 \cite{Eguchi2010} 中：“当我们考虑椭圆亏格时，$K3$ 曲面模空间孤立点处的这些自同构群，是否可能在整个模空间上增强为 $M_{24}$？”其动机在于：椭圆亏格只收到 BPS 态的贡献，因而在整个模空间上不变。于是超对称的保护可以导致一个“meta”对象，它容许我们所要的“meta 对称群”作用。

% [Tutor Note: “meta 群”不是某个已知的有限群，而是一个只在指标层面（椭圆亏格层面）才有意义的对象：模空间不同点的对称群无法同时作用于单个 Hilbert 空间，但可以分别作用在椭圆亏格的各个 BPS 扇区上，再拼成一个更大的代数结构。这正是本节反复出现的“超出微扰”思路的原型。]

在 Kummer 曲面（即由环面轨道流形 $T^4/\ZZ_2$ 得到的 $K3$ 曲面）背景下实现这一想法的具体步骤，见 \cite{Taormina:2011rr,Taormina:2013mda,cheng2015landau}，并在 \cite{GaberdielMatthiasR2017Mmas} 中进一步探索。特别地，三个特定 Kummer 曲面的几何对称性分别由下列群给出：
\begin{equation}
G_0:=\ZZ_2^4\rtimes(\ZZ_2\times\ZZ_2)~,~~G_1:=\ZZ_2^4\rtimes A_4~,~~G_2:=\ZZ_2^4\rtimes S_3~, 
\end{equation}
业已证明可以把它们合并为所谓的八元组群（octad group）$G=\ZZ_2^4\rtimes A_8$。这是 $M_{24}$ 的极大子群，但太大，不能是单个 $K3$ sigma 模型的自同构群。\cite{GaberdielMatthiasR2017Mmas} 的工作提供了证据：$G$ 确实作用在离开模空间中 Kummer 轨迹后仍存活下来的扭结扇区 BPS 态上。这提示：作为 $M_{24}$ 极大子群的 $G$，确实在 Kummer 模空间内充当“统摄性”（overarching）对称性。

% [Tutor Note: 八元组群的名字来自 Steiner 系统 $S(5,8,24)$ 的八元组（octad）：$G=\ZZ_2^4\rtimes A_8$ 正是 $M_{24}$ 中稳定一个八元组的极大子群，阶为 $2^4\cdot 8!=322560$。$G_0,G_1,G_2$ 中的 $A_4$、$S_3$ 都是 $A_8$ 的子群，因而三者都嵌入 $G$；至于把三个 Kummer 模型的对称性拼成 $G$ 这一步，原文依赖 Kummer 轨迹上 BPS 态的显式计算，未展开。]

最后，\cite{UMLG} 在流向红外 $K3$ sigma 模型的 Landau--Ginzburg 模型背景下也探索了这一想法。这一途径的困难之处在于：技术上很难沿模空间中的路径走到一般点——那些点既不对应环面轨道流形，也不具有 Landau--Ginzburg 紫外描述。因此，要把一般 $K3$ sigma 模型中实现的对称性合并起来并不直接。

\myparagraph{低维紧化。} 第二种途径是考虑这样的弦紧化：在模空间的某些给定点上，更大的群作为整个弦论的对称群实现，而不只是微扰弦态的对称群，也不只是 BPS 扇区的对称群。对有 16 个超荷的理论，这只有在非紧维数少于六的紧化中才可能。这一想法最早在 \cite{cheng2010k3} 中讨论，那里考察了 Mathieu 月光对四维 BPS 黑洞的意义——这些黑洞来自紧化在 $K3\times T^2$ 上的 II 型超弦。\cite{Persson:2013xpa} 进一步发展了这一想法，并带来对弦对偶的新洞见 \cite{Persson:2015jka}。

% [Tutor Note: 四维 BPS 黑洞的近视界几何是 $AdS_2\times S^2$，其微观态计数正对应 $K3\times T^2$ 紧化中的 D-膜构型；Mathieu 月光若能组织这些简并度，就提供了“黑洞量子态按散在群分块”的罕见案例。原文用“意义”（implications）而未断言已实现，读时勿夸大。]

最近还证明：紧化到三维的弦论模空间中存在一些点，它们容许 Niemeier 群作为离散对称群 \cite{kachru20173d}。在 IIA 型框架下，这些点由紧化在 $K3\times T^3$ 上给出。这些对称群在理论的 1/2-BPS 态上的作用已有分析 \cite{kachru20173d}；弄清其在 1/4-BPS 态上的作用会很有意思。不过，目前尚不清楚如何在物理设定中取得与本影月光 McKay--Thompson 级数 $\psi^X_g$ 相关的函数。最近 \cite{Zimet:2018dev} 还探索了进一步紧化到 $K3\times T^4$ 的情形。

% [Tutor Note: 三维紧化只有 16 个超荷，1/2-BPS 与 1/4-BPS 的区分相对简单：前者由 $K3$ 上的 BPS 态直接给出，后者才可能带出 $M_{24}$ 的完整作用。原文说“尚不清楚”如何得到 $\psi^X_g$，意指三维设定虽能实现大群，却还没有与月光生成函数对应的指标计算。]

\myparagraph{杂化 $K3$ 紧化。} 这一设定之所以可能相关，是因为其模空间以 ${\cal N}=(4,4)$ $K3$ 共形场论的模空间为子轨迹。因此可以设想：离开这些子轨迹就能看到更多对称性。在技术层面，超对称数目更少，我们计算谱的手段也更少。这条路线在 \cite{cheng2013mathieu,Harrison:2013bya} 中已有一些探索，但仍有大量工作要做。

\myparagraph{五膜（five-brane）动力学。} 本影月光与双标度小弦理论（double scaled little string theory）之间的联系已有研究 \cite{Harvey:2013mda,Harvey:2014cva}；后者描述 IIB 型超弦中 NS5-膜存在时的微扰动力学。

这一想法自然之处在于：它把新的非微扰要素（NS5-膜）纳入物理设定，而双标度小弦理论又有 $ADE$ 分类——这一分类在本影月光的构造中也起重要作用。弄清该设定中存在的有限群对称性会很有意思。

% [Tutor Note: 双标度小弦理论是 NS5-膜背景在特定极限下的可解模型，其谱由 $ADE$ 型仿射代数组织；本影月光的构造也以 ADE 根系（Niemeier 根系统）为骨架，这正是作者说“自然”的原因。但该设定目前只给出谱的定性对应，尚未构造出 $\psi^X_g$ 的物理解释。]

\myparagraph{与某些 VOA 的关系。} 如前几个小节所见，本影月光与 Conway 月光彼此相关，因为二者都与 $K3$ sigma 模型密切相关。\cite{Cheng:2014owa} 表明，Conway 月光模的一个变体带有 $\Co_0$ 的各种保持两平面的子群的作用，其中包括 $M_{23}$，并作为 twining 亏格给出权 0、指标 2 的弱 Jacobi 形式之集。此外，承载 $M_{23}$ 表示的模拟模形式，似乎与 $M_{24}$ 月光中起作用的模拟模形式关系非常密切。更深入地理解这一表面上的联系，并看看能否以某种方式“微调”（tweak）人们更熟悉的 Conway 模以容纳本影模，或许是富有成果的。

% [Tutor Note: 这是五种途径中最“代数”的一条：Conway 模的变体已经能带 $M_{23}$ 这类两平面保持子群，若能把指标从 2 调到本影所需的指标，就可能给出本影模的 VOA 构造。原文用“tweak”一词，说明目前只是方向性设想，没有定理。]

\myparagraph{全然未知。} 由于以上途径迄今都未给出确定答案，一种可能的情形是：我们观察到的 $K3$ 椭圆亏格与本影月光之间的联系只是巧合，而本影月光的物理背景（若存在）完全在别处。

寻找能自然、统一地构造本影月光模的正确物理设定，是当前活跃的研究领域，也是本影月光研究目前的圣杯。

% [Tutor Note: 收尾这句是全文对第 8 节的定调：本影月光的物理起源仍是开放问题；读者若只记住一个结论，应是“K3 椭圆亏格→本影月光”的数学关系已牢固、“物理起源”尚未确立。]
